\section{Scene averaging and contrast magnitude: extended results}\label{app:extended-magnitude}
We compare fit scene by scene, after averaging scenes, and after held-out
rescaling of painter contrasts (Table~\ref{tab:review-calibration}).

\emph{Averaging scenes.} Nano Banana 2's mean painter contrasts have lower error than FLUX's
($.723$ versus $.761$). Evaluating each scene adds a larger variation component
for Nano Banana 2 ($.387$ versus $.040$), reversing that order.
$D$ penalizes variation around the pooled historical target; averaging can
cancel scene-dependent differences.

\begin{table*}[!htbp]
\caption{Retrospective error decomposition and calibration against full-frame references. $D=D_{\rm aggregate}+V_{\rm scene}$. Calibration rescales centered generated contrasts using the other 13 scenes and evaluates them on the omitted scene; it does not produce new images. Bold marks each error column's minimum, not significance.}
\label{tab:review-calibration}
\centering\small
\begin{tabularx}{\linewidth}{@{}l*{3}{>{\centering\arraybackslash}X}@{\hspace{2em}}>{\centering\arraybackslash}X@{}}\toprule
Model & \shortstack{Uncalibrated\\scene error\\$D$} & \shortstack{Error after\\averaging scenes\\$D_{\rm aggregate}$} & \shortstack{Scene\\variation\\$V_{\rm scene}$} & \shortstack{Calibrated\\scene error\\$D_{\rm held}$}\\
\cmidrule(r){1-4}\cmidrule(l){5-5}
GPT Image 1 & 1.572 & 1.239 & 0.333 & 0.645\\
GPT Image 2 & 1.226 & 1.044 & 0.182 & \textbf{0.554}\\
2.5 Flare & 1.738 & 1.483 & 0.255 & 0.741\\
2.5 Sunburst & 1.641 & 1.337 & 0.305 & 0.706\\
Nano Banana 2 & 1.109 & \textbf{0.723} & 0.387 & 0.832\\
FLUX.2 Max & \textbf{0.801} & 0.761 & 0.040 & 0.714\\
\bottomrule\end{tabularx}

\end{table*}


\emph{Rescaling contrasts.} GPT Image 2's aligned response is near one, but its
squared contrast size exceeds twice the reference's ($Q=2.223$, where 1 matches the reference size;
Appendix~\ref{app:magnitude-coverage}).
Shrinking weakens its reference-aligned component, but reduces off-target
differences enough to lower the total error estimate. GPT Image 2's fitted multipliers
scale every contrast coordinate by about 0.45, using only the other 13 scenes each time.

After calibration, GPT Image 2 has the lowest error estimate, $.554$,
compared with FLUX's $.714$; it has lower scores in 12 of 14 held-out folds.
Deleting any scene and refitting the entire procedure preserves the mean
advantage (Appendix~\ref{app:magnitude-coverage}).
The folds overlap, so their win count does not establish statistical superiority.


\section{Source correction: extended comparison}\label{app:extended-source}
\begin{table*}[!htbp]
\caption{Uncalibrated error differences before and after source correction, using audited painting regions and refitted development scaling. Positive differences favor FLUX.2 Max. Full-frame intervals use the 95\% simultaneous procedure across 21 endpoints; corrected intervals reuse that procedure retrospectively and add no confirmatory tests.}
\label{tab:source-comparison}
\centering\small
\begin{tabularx}{\linewidth}{@{}Xcc@{}}\toprule
Model comparison & \shortstack{Full-frame difference\\{[adjusted interval]}} & \shortstack{Crop + scaling correction\\Difference [adjusted interval]}\\\midrule
GPT Image 1 $-$ FLUX & $0.771\;[-0.037,1.580]$ & $0.758\;[0.026,1.491]$\\[2pt]
2.5 Flare $-$ FLUX & $0.937\;[0.256,1.618]$ & $0.864\;[0.193,1.535]$\\[2pt]
2.5 Sunburst $-$ FLUX & $0.840\;[0.058,1.623]$ & $0.705\;[-0.026,1.436]$\\[2pt]
\bottomrule\end{tabularx}

\end{table*}
